\chapter{Joint inference with one-reading analog calibration}
\label{ch:pa-joint-analog}

A bounded analog detector can supply its calibration coordinate with a single
reading. No unchanged-state twin is needed for the conservative construction
below. Combining that coordinate with the two dynamical scores gives a single
confidence bound. The direct semigroup dissipation inequality then converts
the same covered contrast into a substantially stronger physical certificate.
The statistical construction extends the joint-calibration and direct-analog
chapters; its ingredients remain the reversible quadratic moment inequality,
conditional Jensen, and Hoeffding concentration.

\section{The observable second moment has the useful inequality sign}
Let $Z\in[-1,1]$ be the output of a fixed noninvasive memoryless detector, with
conditional mean $f(x)=\mathbb E[Z\mid X=x]$. Use the stationary notation
\[
 A=\mathbb E_\pi f^2,\qquad B=\mathbb E_\pi Z^2,
 \qquad C_k=\langle f,P^k f\rangle_\pi,
 \qquad P=e^{\Delta Q}.
\]
Conditional detector noise at the two endpoints of a positive-lag pair is
independent. Consequently
\[
 B=A+\mathbb E_\pi\operatorname{Var}(Z\mid X)\geq A.
\]
\[
 W_0=\frac{1+Z^2}{2}\in[1/2,1],\qquad
 W_k=\frac{1+Z_0Z_{k\Delta}}2\in[0,1]\quad(k=1,2),
\]
and their means are $p_0=(1+B)/2$ and $p_k=(1+C_k)/2$.
Here $C_0=A$; the observable squared reading estimates $B$, not $A$.

\begin{proposition}[One-reading analog reversible support]
For fixed $0\leq s\leq1$, let $a_s=(-s^2,2s,-1)$ and
$b_s=-(1-s)^2/2$. Every reversible model satisfying this detector contract obeys
\begin{equation}
 a_s\cdot p-b_s
 =\frac{-s^2B+2sC_1-C_2}{2}
 =-\frac{\|(P-sI)f\|_\pi^2+s^2(B-A)}2\leq0.
 \label{eq:pa-analog-support}
\end{equation}
The support $b_s$ is sharp across the allowed detector class.
\end{proposition}
\begin{proof}
For reversible $Q$, the operator $P$ is self-adjoint, so
$\|(P-sI)f\|_\pi^2=C_2-2sC_1+s^2A$.
The additional conditional-variance term is nonnegative. The deterministic
zero-output detector attains equality for every $s$.
\end{proof}

The detector's noise costs population signal:
\[
 g_s^{\rm analog}=g_s^{\rm ideal\ twins}-\frac{s^2}{2}(B-A).
\]
Avoiding twins therefore does not guarantee better power. When the detector
noise dominates its state-dependent signal, the conservative square may erase
the violation entirely. For a state-independent fair detector with outputs
$\pm1$, $A=0$, $B=1$, and the analog gap is $-s^2/2$, even though its true
conditional mean contains no state information. This example also shows why
one cannot silently reinterpret a squared noisy reading as $A$.

\section{One confidence event, with every acquisition counted}
Take $n_0$ independent stationary single readings and $n_1,n_2$ independent
stationary pairs at the two lags. Let $\widehat p_k$ denote their score means.
For a fixed or independently piloted $s$, define
\begin{align}
 r_s&=\sqrt{\frac{\log(1/\alpha)}2
       \left(\frac{s^4}{4n_0}+\frac{4s^2}{n_1}+\frac1{n_2}\right)},
       \label{eq:pa-analog-radius}\\
 \underline g_s&=\left[
   -s^2\widehat p_0+2s\widehat p_1-\widehat p_2-b_s-r_s
   -\sum_{k=0}^2|a_{s,k}|(e_k+\delta_k)
 \right]_+.
 \label{eq:pa-analog-gap}
\end{align}
Here $e_k$ bounds numerical error in an empirical mean, and $\delta_k$ bounds
the conditional mean discrepancy from its ideal target given the entire past.
Independent stationary acquisitions have $\delta_k=0$. The conditional version
requires a valid external bound, not merely marginal stationarity. Hoeffding's
bound, or its conditional exponential-moment proof, gives the stated event
with failure at most $\alpha$. Independently established contract failure can
be added by a union bound. The API records these premises without verifying
them from the score summaries.

The factor $1/4$ in the first variance term is the square of the score's range
length $1/2$. It does not eliminate that term. For example, use a reversible
reset generator with stationary distribution $(1/2,1/4,1/4)$, detector values
$(0,1,-1)$, decay $1/2$, and $s=1/2$. Then $A=1/2$ and the ideal score means
are $(3/4,5/8,9/16)$, exactly on the support. One calibration reading is zero
with probability $1/2$, giving $\widehat p_0=1/2$. Even if the two dynamical
means were known exactly, treating that calibration reading as exact would
create a false gap $1/16$. Equation~\eqref{eq:pa-analog-radius} retains the
uncertainty and prevents this example from giving a positive certificate.

This design uses
\[
 N_{\rm acquisitions}=n_0+n_1+n_2,\qquad
 N_{\rm readouts}=n_0+2n_1+2n_2,\qquad
 T_{\rm evolution}=\Delta n_1+2\Delta n_2.
\]
Preparation, reset, sensor and calibration overhead must also be counted when
known. None is assigned zero physical energy here. The exact square-summary
helper rounds the score with a declared numerical radius; it bounds the
rounding error of stored digital readings, not unknown sensor calibration
error.

\section{A direct physical conversion on the same event}
Under the stationary even-variable Markov assumptions of the direct spectral
dissipation theorem, the detector mean has range at most two and
\[
 -s^2B+2sC_1-C_2\leq -s^2A+2sC_1-C_2.
\]
Thus the covered gap in~\eqref{eq:pa-analog-gap} supplies
\begin{equation}
 \sigma\geq\frac{4\underline g_s}{\Delta},\qquad
 \sigma\geq\frac{\phi(2\underline g_s)}{\Delta},\qquad
 \phi(u)=2\sqrt u\,\operatorname{atanh}\sqrt u,
 \quad 0\leq u<1.
 \label{eq:pa-analog-direct-bridge}
\end{equation}
The implementation encloses $\phi$ with a positive rational series. The
interval is an enclosure of this lower-bound formula, not an interval for
the true entropy production. An incompatible normalized defect is flagged
on the coverage event; it is not silently clipped or presented as empirical
falsification.

For comparison, the earlier additive-reversible record argument gives
\begin{equation}
 \sigma\geq
 \frac{8\underline g_s^2}{(2s\sqrt\Delta+\sqrt{2\Delta})^2}.
 \label{eq:pa-analog-old-bridge}
\end{equation}
The zero-lag squared-reading law is unchanged by additive reversibilization
because the stationary law and detector channel are unchanged. It therefore
contributes sampling uncertainty and apparatus expenditure but no comparison
divergence in~\eqref{eq:pa-analog-old-bridge}. The stronger direct argument
uses the semigroup identity and forward/reverse comparison instead. Both
physical conversions reuse the same confidence event; they spend no additional
$\alpha$. Different choices of $s$ still require independent selection or a
simultaneously valid procedure.

\section{Frozen comparison on the same analog data}
The benchmark fixes all models, parameters, counts, rules and random seed
before drawing its 1,600 held-out synthetic datasets. Each of four cases has
400 independent datasets. A 20-state circulant CTMC is observed through a
cosine detector rounded to denominator 4096. Each dataset contains 4,000
single-reading calibration acquisitions and 3,000 pairs at each of lags
$1/10$ and $1/5$. Every compared method therefore receives the same
10,000 acquisitions, 16,000 detector readings and 900 declared units of
evolution time.
These comparisons match acquisition resources. Their analysis-computation
time and energy are unmeasured and are not asserted equal.

The base contrast has $s=1/2$, reproducing the earlier difference of positive-lag
scores. A second contrast fixes $s=3/4$ before production; ``tuned'' refers to
this predetermined design, not a fit to the held-out records. The split methods
call the sixth edition's actual analog calibration and score APIs, with
$\alpha/2=1/40$ for each stage. The joint methods use $\alpha=1/20$. Virtual
coin methods retain a draw from $\operatorname{Bernoulli}(W_k)$ for each same
analog pair. The calibration singles are common to every method.

\begin{center}
\begin{tabular}{lrrr}
\toprule
Method & Driven & Equilibrium & Noisy driven\\
\midrule
Joint analog, $s=1/2$ &399/400&0/400&0/400\\
Joint analog, $s=3/4$ &399/400&0/400&0/400\\
Joint coin, $s=3/4$ &361/400&1/400&0/400\\
Split analog, $s=1/2$ &377/400&0/400&0/400\\
Split analog, $s=3/4$ &379/400&0/400&0/400\\
Split coin, $s=1/2$ &293/400&0/400&0/400\\
\bottomrule
\end{tabular}
\end{center}

The driven rates are 20 forward and 1 backward; equilibrium rates are both
$21/2$. The noisy detector outputs half the rounded cosine plus an independent
fair noise of amplitude $1/2$. Its squared-reading calibration is too
conservative here, and every method returns zero in all 400 datasets.
An exact population audit explains this failure. For $B>0$, completing the
square gives
\[
 \sup_{s\in\mathbb R}g_s^{\rm analog}
 =\frac12\left(\frac{C_1^2}{B}-C_2\right).
\]
Rational uniformization of this fixed noisy ring encloses both correlations
and certifies that the right side is below $-0.007$. Its numerical value is
approximately $-0.00789686$. The entire two-lag quadratic family therefore
lacks a positive population gap here. More samples at these same settings
cannot repair this obstruction; different lags, additional noise calibration
or different observables remain possible. This retrospective exact audit does
not change the frozen inference design or any reported decisions.
The single equilibrium rejection by the joint coin method is reported rather
than removed; the methods have individual 5\% failure allowances.

Keeping the older quadratic physical conversion isolates the statistical
improvement. At $s=1/2$, the driven mean lower bound rises from $0.00145847$
for split analog to $0.00351542$ for joint analog. At $s=3/4$, it rises from
$0.00164098$ to $0.00494257$. Applying the new direct bridge to that same
joint $s=3/4$ event gives a mean linear bound $0.870376$ and a mean nonlinear
bound $0.884824$: the latter is about 179 times the older quadratic mean.
The simulated model's actual rate is $56.9189$, so even the improved mean is
only 1.55\% of the microscopic value. This large remaining gap is explicit.

The virtual coin's greater variance can increase the average of a clipped
quadratic lower bound even while reducing detection probability. Indeed its
joint $s=3/4$ quadratic mean is $0.00639217$, above the analog value, despite
361 rather than 399 detections. This is why an average positive-part bound
alone does not establish estimator quality or universal statistical dominance.
The matched table reports detection counts, means, medians and Monte Carlo
intervals in the accompanying results.

\section{A detector change remains a physical failure}
The fourth case uses an equilibrium ring with rates 70 in both directions,
but replaces the calibration detector by zero gain while retaining the cosine
science detector. This deliberately violates the common-channel contract.
Both joint analog methods and both split analog methods falsely return positive
values in all 400 datasets. The split coin returns 395 positives and the joint
coin 400. Correcting the zero-lag conditional-mean discrepancy by the general
bound $\delta_0=1/2$ makes the reported repair return zero in all datasets.
This repair is deliberately conservative; obtaining a useful smaller bound
requires actual detector calibration.

The archive retains every dataset's exact integer sufficient statistics and
the first raw analog bank in each case. Rational score means can be replayed
without drawing new coins or trajectories. The six methods are separate
predeclared analyses; their unadjusted maximum is not advertised as a 95\%
certificate. These are reproducible simulations, not new measured physical
data or a universal noise-removal procedure. Physical use requires stationary
preparation, the appropriate reversal convention, detector independence and
memory control, time calibration, and a justified thermodynamic embedding for
conversion from entropy rate to dissipated power.
